\documentclass[11pt,a4paper]{article}
\usepackage{planos-ufmt}
\begin{document}
\CapaPlano{Teorema de Bayes}{UFMT Várzea Grande, área Matemática Aplicada/Probabilidade, prova didática 2025}{Probabilidade I; graduação com eventos, condicionamento e partições.}
\section{Objetivos e estratégia}
Provar Bayes a partir de probabilidade condicional e probabilidade total, calcular uma posterior e explicar a armadilha da taxa-base. Metodologia: problema inicial com duas hipóteses, prova algébrica e análise interpretativa.
\Cronograma{Pergunta: um teste positivo implica alta probabilidade do evento investigado?}{Definir $P(A\mid B)$, hipóteses de positividade e partição por eventos $H_i$.}{Provar a lei da probabilidade total e o teorema de Bayes com denominador correto.}{Resolver exemplo de prevalência de $1\%$, sensibilidade $90\%$ e falso positivo $5\%$; interpretar $15{,}38\%$.}{Fechar contrapondo $P(T^+\mid H)$ a $P(H\mid T^+)$.}
\section{Demonstrações rigorosas}
Se $P(B)>0$, define-se $P(A\mid B)=P(A\cap B)/P(B)$. Se $\{H_1,\dots,H_n\}$ é partição do espaço amostral em eventos de probabilidade positiva, então os eventos $B\cap H_i$ são disjuntos e têm união $B$. Portanto
\[
P(B)=\sum_{i=1}^nP(B\cap H_i)
=\sum_{i=1}^nP(B\mid H_i)P(H_i).
\]
\begin{teorema}[Bayes]
Se $P(B)>0$ e $\{H_1,\dots,H_n\}$ é partição com $P(H_i)>0$, então
\[
P(H_j\mid B)=
\frac{P(B\mid H_j)P(H_j)}
{\sum_{i=1}^{n}P(B\mid H_i)P(H_i)}.
\]
\end{teorema}
\begin{proof}
Pela definição de probabilidade condicional,
$P(H_j\mid B)=P(B\cap H_j)/P(B)$.
Também $P(B\cap H_j)=P(B\mid H_j)P(H_j)$, e a lei da probabilidade total fornece o denominador. Substituindo, obtemos a identidade.
\end{proof}
\section{Exemplo com resposta e interpretação}
Seja $H$ o evento de uma condição presente em $1\%$ da população, e $T^+$ o teste positivo. Suponha $P(T^+\mid H)=0{,}90$ e $P(T^+\mid H^c)=0{,}05$. Logo
\[
P(H\mid T^+)=
\frac{0{,}90\cdot0{,}01}
{0{,}90\cdot0{,}01+0{,}05\cdot0{,}99}
=\frac{0{,}009}{0{,}0585}\approx0{,}1538.
\]
A posterior é aproximadamente $15{,}38\%$, não $90\%$. Para explicar no quadro: em dez mil casos hipotéticos, cerca de $100$ teriam a condição, com $90$ positivos verdadeiros; entre os outros $9900$, haveria $495$ falsos positivos, produzindo $90/(90+495)=15{,}38\%$. \textbf{Limite:} frequências e taxas são pressupostos do exemplo, não dados empíricos gerais.
\Fechamento{Esquerda: partição; centro: equação de Bayes; direita: tabela de frequências entre 10 mil casos.}{Por que $P(B\mid A)$ e $P(A\mid B)$ podem ser muito diferentes? Porque seus denominadores representam populações de referência diferentes.}
\section{Autoavaliação específica}
\begin{itemize}
\item Declarei $P(B)>0$ e eventos condicionantes de probabilidade positiva?
\item Provei a decomposição do denominador por uma partição?
\item Expliquei a discrepância entre sensibilidade e valor preditivo positivo?
\end{itemize}
\RubricaUFMT
\section{Referências}
M. N. Magalhães e A. C. P. Lima, \emph{Noções de Probabilidade e Estatística}; bibliografia de Probabilidade do edital de referência.

\end{document}
